\subsection{Graphs of Whole Number Powers} 
Consider these example graphs.
\begin{center}
\begin{tikzpicture}[x=0.5\linespace,y=0.5\linespace]
\setgraphsmall
\useasboundingbox (.75\textwidth,-\value{graphsize}) rectangle (1.0\textwidth,\value{graphsize});
\foreach \sh in {2,...,5}
	{
	\begin{scope}[xshift=\sh*0.25*\textwidth]
		\GraphSmall
		\draw[domain={-\value{graphsize}^(1/\sh)}:{\value{graphsize}^(1/\sh)},samples=100,style=graph] 
		      plot (\x,{(\x)^(\sh)});
		\draw (\value{graphsize}+.5,-\value{graphsize}-.5) node[anchor=south east]{\small$y=x^\sh$};
	\end{scope}
	}
\end{tikzpicture}
\end{center}
\begin{itemize}
\item If $n$ is \makeblank[1in], the graph is a \makeblanksmall-shaped curve that passes through \blankpair, \blankpair, and \blankpair.
\item If $n$ is \makeblank[1in], the graph is a \makeblanksmall-shaped curve that passes through \blankpair, \blankpair, and \blankpair.
\item The \makeblank[1in] the value of $n$, the \makeblank[1.5in] the way the graph turns.
\end{itemize}
To graph $f(x)=a(x-h)^n+k$, we keep the general shape with a turn at \blankpair\linebreak and use two points on either side. If $a$ is a \makeblank[1in] number, the graph is\linebreak \makeblank.
\begin{example}
Graph 
\[
f(x)=-\frac{1}{2}\left(x-2\right)^3+1
\]
\begin{lpic}{}{10}
	\regtext{2}{-2}{$x$};
	\regtext{4}{-2}{$y$};
	\draw (3,-1) -- (3,-7);
	\draw (1,-2) -- (5,-2);
	\draw[dashed] (1,-4) rectangle (5,-5);
\end{lpic}
\end{example}